\section{Discussion}
The current evidence favors adapting the classifier to a bounded receiver support bank over training an attentive context-fusion module for this task. T3A's rule is simple, but its success is not unconditional: the smallest tested banks degrade recognition, and probability calibration changes differently from class accuracy. A deployment that cannot supply sufficiently informative support may prefer a source-only decision rule. The diagnostic does not establish an operationally optimal bank size.

The P2 result constrains interpretation. Receiver-matched support changes predictions, yet the net gain relative to P0 is small and heterogeneous. Better results than corrupted-context controls show mechanism sensitivity, not superiority to a competent baseline. Neither attention weights nor oracle composition experiments establish causality. We retain the neutral outcome as a negative architectural finding within the evaluated setup.

The study's most defensible methodological value is information accounting. Source-only receiver invariance, prototype adjustment, input-statistic normalization, context conditioning, and supervised support fine-tuning solve different information-access problems. Receiver holdout alone does not make them equivalent. Reporting support size, query isolation, parameter updates, and label availability is necessary for fair comparison.

Evidence timing matters. Methods remediation and the subsequent benchmark use the same WiSig subset; the later statistical family was selected with earlier outcomes known. The strong T3A pattern is a reason to replicate, not a substitute for replication. An independent dataset must verify physical receiver identities, lawful provenance, compatible representation, and target-neutral support construction before new results. Prospective collection should acquire calibration and query sessions separately and preserve clock and session provenance rather than reconstructing episodes from class-organized files.

Software readiness and scientific readiness differ. A tested Shen adapter cannot resolve licensing or prove that a centered short I/Q crop preserves LoRa semantics. A crash-tested collection journal cannot verify an SDR clock or guarantee mixed transmitter activity. These are actionable gaps, not additional evidence for the current claim.
